\section{Introduction}\label{sec:intro}

Machine-learning network intrusion detectors are mostly evaluated on public benchmarks, where many report macro-F1 close to the maximum. Such scores say little about an IoT gateway, which has a few CPU cores, sees traffic in time order, builds flows itself and must act on its verdicts. The gap between benchmark and operation is long known~\cite{sommer2010outside} and well documented for security machine learning~\cite{arp2022dos,apruzzese2023sok}; a recent study reports a random forest with hold-out F1 0.97 whose attack recall fell to 0.0095 in operational validation on a labelled laboratory capture~\cite{tang2026hybrid}.

The sources of this gap are usually studied one at a time: label errors~\cite{engelen2021troubleshooting,liu2022error,flood2024smells}, splits that ignore time~\cite{pendlebury2019tesseract,arp2022dos}, and flow exporters that cut flows differently~\cite{pinto2025flowexporter,vitorino2025revisiting}. Live prototypes rarely check that they classify the flows the model was evaluated on, and automatic blocking is rarely scored by the hosts it blocks wrongly.

We audit one ordinary detector along the whole path: a decision tree over 54 NFStream~\cite{aouini2022nfstream} flow statistics, trained on ten captures of CIC-IoT2023~\cite{neto2023ciciot2023}, with the offline steps repeated on ten IoT-23 scenarios~\cite{garcia2020iot23}. We propose no new detection method. At each of four steps (labels, splits, live gateway, response) we ask what the previous step hid and which check reveals it. The contributions are:
\begin{itemize}
\item \emph{The cost of file-name labels, on two datasets} (Section~\ref{sec:labels}). A tree trained on file-name labels flags 71.8\% of benign CIC-IoT2023 flows (against our evidence labels) and, in a robustness test against the IoT-23 authors' Zeek labels, 97--100\% of background flows in every malware capture. The CSVs released with CIC-IoT2023 allow only file-name labels.
\item \emph{A live-fidelity check and the failure modes it exposed} (Section~\ref{sec:live}): a forwarding pitfall, idle flow meters and an NFStream export-ordering defect, reported as nfstream\#262~\cite{nfstream262}, with a fix submitted as PR \#263~\cite{nfstream263}.
\item \emph{Response scored in false blocks} (Section~\ref{sec:response}), with the trade-off over the number of verdicts required per host pair.
\item \emph{A reusable four-check audit} (Section~\ref{sec:checklist}), implemented as scripts in the repository and run on two datasets.
\end{itemize}
